\documentclass[12pt]{article}

% ---------- Page setup: 1-inch margins, 12-pt font ----------
\usepackage[margin=1in]{geometry}
\usepackage[T1]{fontenc}
\usepackage{lmodern}
\usepackage{microtype}
\usepackage{amsmath,amssymb}
\usepackage{xcolor}
\usepackage{enumitem}
\usepackage[compact]{titlesec}
\usepackage[authoryear,round]{natbib}
\usepackage[colorlinks=true,linkcolor=blue!50!black,citecolor=blue!50!black,urlcolor=blue!50!black]{hyperref}

% Compact spacing so the body fits in two pages
\titleformat{\section}{\large\bfseries}{\thesection}{0.5em}{}
\titlespacing*{\section}{0pt}{6pt}{2pt}
\setlist{nosep,leftmargin=*}
\setlength{\parindent}{0pt}
\setlength{\parskip}{3pt}

\begin{document}

% ---------- Title block (manual, to save space) ----------
\begin{center}
{\Large\bfseries Does Directional Dropout Make Refusal Harder to Remove?}\\[1pt]
{\large Self-Distillation under Random Ablation of the Refusal Direction}\\[3pt]
Ryan Xu \quad Julien Schweighofer \quad Michelle Koga \quad Rushil Saraf\\
{\small Topic: Safety of open-weight models}
\end{center}
\vspace{-6pt}

% ---------- Abstract (5 sentences) ----------
\textbf{Abstract.}
Anyone can remove the safety training of an open-weight model.
The most common method, abliteration, finds one direction in the model's activations that controls refusal and deletes it from the weights.
We test a defense we call directional dropout: during training we randomly delete the refusal direction, and the model has to keep refusing anyway.
We train the model to copy the outputs of the original model (self-distillation), so it learns a new way to represent refusal rather than new behavior.
We then attack the trained Qwen3.5-0.8B with abliteration, pruning, prefilling and LoRA.

% ============================================================
\section{Introduction and Background}
\textbf{Safety is easy to remove.}
\citet{arditi2024refusal} find that refusal is controlled by one direction in the residual stream, and that deleting it from the weights (abliteration) stops the model from refusing.
\citet{wollschlager2025geometry} show that refusal can also be carried by several independent directions.
\citet{wei2024brittleness} show that pruning about 3\% of the weights breaks safety while the model stays useful.
\citet{qi2024finetuning} show that fine-tuning on a handful of harmful examples undoes safety training.
\citet{qi2025shallow} show that safety training mostly affects the first few tokens of a response, so forcing a compliant start (prefilling) bypasses it.

\textbf{Existing defenses.}
TAR \citep{tamirisa2025tar} uses adversarial meta-learning to resist harmful fine-tuning, but \citet{kuo2026simple} show that abliteration and prefilling still break it.
ReFAT \citep{yu2025refat} and DeepRefusal \citep{xie2025deeprefusal} delete the refusal direction during safety training, which makes models harder to abliterate.
Both also train on new refusal examples, so it is unclear whether the gain comes from more safety training or from a new representation of refusal.

\textbf{Dropout.}
Dropout randomly turns off units during training, so the network cannot rely on any one of them \citep{srivastava2014dropout}.
\citet{liu2025dropdropout} find that models pretrained without dropout are easier to edit, which suggests that dropout stores knowledge in more than one place.

% ============================================================
\section{Key Questions}
\begin{enumerate}
  \item Does training a model to refuse while its refusal direction is randomly deleted make it harder to abliterate?
  \item Does it have to be the refusal direction, or would a random direction or noise work as well?
  \item How many directions must an attacker then remove, and does the defense help against pruning, prefilling or LoRA?
\end{enumerate}

% ============================================================
\section{Hypothesis}
We expect the model to learn to store refusal in several directions, so that removing one direction no longer breaks it.
We expect the random-direction and noise controls to do little.
We expect recomputing the direction during training to matter, since with a fixed direction the model can simply move refusal to one new direction.
We do not expect the defense to help against pruning or LoRA \citep{qi2024finetuning,kuo2026simple}.

% ============================================================
\section{Methods}
\textbf{Self-distillation.}
A frozen copy of the safety-tuned model is the teacher, and a second copy is the student.
During training we perturb the student and minimize the KL divergence between its next-token distribution and the teacher's, on the teacher's own responses \citep{hinton2015distilling}.
The student starts as an exact copy, so it only changes to make up for the perturbation, and its behavior should stay the same.

\textbf{Directional dropout.}
The refusal direction is the difference between the average activations on harmful and harmless prompts \citep{arditi2024refusal}.
At each layer, with probability 0.25 or 0.5, we remove it from the student's activations.
We recompute it from the student every few hundred steps.

\textbf{Conditions.}
(1) The original model; (2) random noise of the same size; (3) deleting a random direction; (4) directional dropout; (5) directional dropout with a fixed direction; (6) directional dropout on the top three directions.

\textbf{Models, code and data.}
We use Qwen3.5-0.8B with thinking off, and Llama~3.2~1B~Instruct if time allows.
We build on the code of \citet{arditi2024refusal} and \citet{wei2024brittleness}, with PEFT for LoRA and vLLM for generation.
We train on about 5{,}000 benign and 1{,}000 harmful prompts from AdvBench \citep{zou2023universal} and HarmBench \citep{mazeika2024harmbench}.
The targets are the teacher's own responses, so we write no new harmful content.
We test on held-out prompts from HarmBench, StrongREJECT \citep{souly2024strongreject} and XSTest \citep{rottger2024xstest}.

\textbf{Attacks and evaluation.}
We attack each trained model with abliteration of 1 to 10 directions, pruning of the most safety-relevant weights, prefilling, and a few LoRA steps.
StrongREJECT's judge scores harmfulness; perplexity, MMLU and XSTest measure usefulness and over-refusal.
We report training and validation loss curves and 3 seeds per condition with 95\% confidence intervals, and release only defended checkpoints.

\textbf{Compute.}
Each training run takes under an hour on an L40S, so training costs about 25 L40S-hours.
Evaluating 16 models takes about 50 L4-hours; the Llama replication would roughly double both.

% ============================================================
\begingroup
\small
\setlength{\bibsep}{1pt}
\begin{thebibliography}{99}

\bibitem[Arditi et~al.(2024)]{arditi2024refusal}
A.~Arditi, O.~Obeso, A.~Syed, D.~Paleka, N.~Panickssery, W.~Gurnee, and N.~Nanda.
\newblock Refusal in language models is mediated by a single direction.
\newblock In \emph{NeurIPS}, 2024.

\bibitem[Hinton et~al.(2015)]{hinton2015distilling}
G.~Hinton, O.~Vinyals, and J.~Dean.
\newblock Distilling the knowledge in a neural network.
\newblock \emph{arXiv:1503.02531}, 2015.

\bibitem[Kuo et~al.(2026)]{kuo2026simple}
K.~Kuo, V.~Smith, and C.~Yadav.
\newblock Open-weight LLM fine-tuning defenses are susceptible to simple attacks.
\newblock \emph{arXiv:2605.26526}, 2026.

\bibitem[Liu et~al.(2025)]{liu2025dropdropout}
H.~Liu, J.~Bauer, and C.~D. Manning.
\newblock Drop dropout on single-epoch language model pretraining.
\newblock In \emph{Findings of ACL}, 2025.

\bibitem[Mazeika et~al.(2024)]{mazeika2024harmbench}
M.~Mazeika, L.~Phan, X.~Yin, A.~Zou, Z.~Wang, N.~Mu, E.~Sakhaee, N.~Li, S.~Basart, B.~Li, D.~Forsyth, and D.~Hendrycks.
\newblock HarmBench: A standardized evaluation framework for automated red teaming and robust refusal.
\newblock In \emph{ICML}, 2024.

\bibitem[Qi et~al.(2024)]{qi2024finetuning}
X.~Qi, Y.~Zeng, T.~Xie, P.-Y. Chen, R.~Jia, P.~Mittal, and P.~Henderson.
\newblock Fine-tuning aligned language models compromises safety, even when users do not intend to!
\newblock In \emph{ICLR}, 2024.

\bibitem[Qi et~al.(2025)]{qi2025shallow}
X.~Qi, A.~Panda, K.~Lyu, X.~Ma, S.~Roy, A.~Beirami, P.~Mittal, and P.~Henderson.
\newblock Safety alignment should be made more than just a few tokens deep.
\newblock In \emph{ICLR}, 2025.

\bibitem[R{\"o}ttger et~al.(2024)]{rottger2024xstest}
P.~R{\"o}ttger, H.~R. Kirk, B.~Vidgen, G.~Attanasio, F.~Bianchi, and D.~Hovy.
\newblock XSTest: A test suite for identifying exaggerated safety behaviours in large language models.
\newblock In \emph{NAACL}, 2024.

\bibitem[Souly et~al.(2024)]{souly2024strongreject}
A.~Souly, Q.~Lu, D.~Bowen, T.~Trinh, E.~Hsieh, S.~Pandey, P.~Abbeel, J.~Svegliato, S.~Emmons, O.~Watkins, and S.~Toyer.
\newblock A StrongREJECT for empty jailbreaks.
\newblock In \emph{NeurIPS Datasets and Benchmarks Track}, 2024.

\bibitem[Srivastava et~al.(2014)]{srivastava2014dropout}
N.~Srivastava, G.~Hinton, A.~Krizhevsky, I.~Sutskever, and R.~Salakhutdinov.
\newblock Dropout: A simple way to prevent neural networks from overfitting.
\newblock \emph{JMLR}, 15:1929--1958, 2014.

\bibitem[Tamirisa et~al.(2025)]{tamirisa2025tar}
R.~Tamirisa, B.~Bharathi, L.~Phan, A.~Zhou, A.~Gatti, T.~Suresh, M.~Lin, J.~Wang, R.~Wang, R.~Arel, A.~Zou, D.~Song, B.~Li, D.~Hendrycks, and M.~Mazeika.
\newblock Tamper-resistant safeguards for open-weight LLMs.
\newblock In \emph{ICLR}, 2025.

\bibitem[Wei et~al.(2024)]{wei2024brittleness}
B.~Wei, K.~Huang, Y.~Huang, T.~Xie, X.~Qi, M.~Xia, P.~Mittal, M.~Wang, and P.~Henderson.
\newblock Assessing the brittleness of safety alignment via pruning and low-rank modifications.
\newblock In \emph{ICML}, 2024.

\bibitem[Wollschl{\"a}ger et~al.(2025)]{wollschlager2025geometry}
T.~Wollschl{\"a}ger, J.~Elstner, S.~Geisler, V.~Cohen-Addad, S.~G{\"u}nnemann, and J.~Gasteiger.
\newblock The geometry of refusal in large language models: Concept cones and representational independence.
\newblock In \emph{ICML}, 2025.

\bibitem[Xie et~al.(2025)]{xie2025deeprefusal}
Y.~Xie, Y.~Zhang, T.~Liu, D.~Ma, and T.~Liu.
\newblock Beyond surface alignment: Rebuilding LLMs safety mechanism via probabilistically ablating refusal direction.
\newblock In \emph{Findings of EMNLP}, 2025.

\bibitem[Yu et~al.(2025)]{yu2025refat}
L.~Yu, V.~Do, K.~Hambardzumyan, and N.~Cancedda.
\newblock Robust LLM safeguarding via refusal feature adversarial training.
\newblock In \emph{ICLR}, 2025.

\bibitem[Zou et~al.(2023)]{zou2023universal}
A.~Zou, Z.~Wang, N.~Carlini, M.~Nasr, J.~Z. Kolter, and M.~Fredrikson.
\newblock Universal and transferable adversarial attacks on aligned language models.
\newblock \emph{arXiv:2307.15043}, 2023.

\end{thebibliography}
\endgroup

% ============================================================
\appendix
\section*{Appendix: Use of Generative AI Tools}
% Does not count toward the page limit (course policy).
We used Claude (Anthropic), through the Claude Code desktop app, to restructure our draft proposal into this format, and to condense it to two pages and to check the bibliographic details (authors, venues, arXiv numbers) of the cited papers against arXiv, ACL Anthology and conference pages.

\end{document}
